\section{Introduction}
\label{sec:introduction}

HTTP security decisions increasingly involve evidence that is difficult to reduce to a single signature. A request may contain a recognizable injection string, a potentially sensitive object identifier, ambiguous authorization context, or natural-language content intended for an AI-enabled application. Signature engines remain fast and auditable, but their coverage is bounded by encoded rules. Learned classifiers can generalize beyond exact signatures, yet their reliability depends on the data distribution, the available context, and the cost assigned to false alarms. Instruction-following language models add another option because they can jointly interpret paths, headers, bodies, and contextual fields while returning a structured security judgment. The same instruction-following behavior also creates risk: attacker-controlled content may attempt to redirect the classifier, and a fluent answer may conceal an unjustified decision~\citep{perez2022promptinject,greshake2023indirect,schulhoff2023hackaprompt}.

Recent cybersecurity research has begun to evaluate local language models as traffic-analysis aids rather than assuming that general language competence transfers directly to attack detection. For example, a recent study in the \emph{Journal of Cyber Security Technology} found that local models could recognize some language-based attack indicators but remained unreliable across packet formats, attack types, and benign traffic~\citep{paukkeri2026local}. That result motivates a more specific question at the HTTP decision boundary: \emph{when a compact model is adapted to a structured request-classification task, which improvements are real, and which deployment-critical failures remain?}

This paper addresses that question through an empirical study of a 4B-parameter model adapted with QLoRA~\citep{hu2021lora,dettmers2023qlora}. The objective is not to introduce a new learning algorithm or model architecture. Instead, the study characterizes the reliability limits of compact language models used for request-side HTTP security decision support. This distinction matters because a statistically significant improvement over an untuned base model does not by itself establish operational usefulness. A model may gain attack recall while producing too many false positives, may satisfy a JSON schema while making the wrong security decision, or may force insufficient-context cases into a binary verdict instead of abstaining.

The evaluation therefore treats leakage control and failure decomposition as part of the cybersecurity evidence. The 413-sample \challenge{} challenge set is external to all training records and few-shot examples. It is excluded from prompt selection, threshold selection, hyperparameter tuning, and model selection. Prediction and scoring are separated: inference reads only model-visible request evidence, writes raw outputs, and terminates before the scoring process loads frozen labels. The study evaluates \texttt{safe}, \texttt{attack}, and \texttt{unknown} as distinct classes and reports safe false-positive rate, attack recall, abstention behavior, schema validity, latency, memory, artifact size, and paired statistical tests. Multiple model artifacts and message-serialization paths are tested because security behavior can change even when nominal weights and generation settings appear equivalent.

The principal result is deliberately mixed. The adapted model improves verdict macro-F1 from 0.4175 for zero-shot Qwen3-4B to 0.4658, with a paired difference of 0.0482 and a 95\% bootstrap confidence interval of 0.0113--0.0835. Attack recall reaches 0.9072. However, the model falsely labels 43 of 209 safe requests as attacks and predicts no unknown verdict for any of 107 insufficient-context cases. A change from separate system/user messages to one concatenated message alters 11 of 413 verdicts. These findings support use as an attack-sensitive shadow signal or analyst triage aid, but they contradict an autonomous-blocking interpretation.

The paper makes four empirical contributions:
\begin{enumerate}
    \item \textbf{A leakage-controlled external evaluation.} The study binds a frozen request-only challenge set, separated inference and scoring processes, content hashes, raw outputs, and paired statistical analysis into one auditable evaluation record.
    \item \textbf{A same-input comparison across security decision systems.} Eight systems are evaluated on identical model-visible requests, including repository rules, a character n-gram classifier, zero-shot and few-shot base-model prompts, the QLoRA adapter, merged artifacts, and GGUF serving paths.
    \item \textbf{A decomposition of cybersecurity failure modes.} The analysis separates attack detection, safe-request false positives, unknown-class abstention, structured-output validity, prompt-injection behavior, quantization effects, and message-serialization sensitivity.
    \item \textbf{Operationally bounded conclusions.} The results identify where compact models may assist triage and where they remain unsafe, particularly when false-positive cost and insufficient context are ignored.
\end{enumerate}

The remainder of the paper reviews related cybersecurity and language-model studies, defines the task and model, documents the frozen data and evaluation protocol, presents comparative and ablation results, and discusses the implications for security decision support.
